\chapter{تکرار با حلقه‌ها}%
\label{ch:loops}

به سومین و آخرین آجرِ سازندهٔ الگوریتم‌ها رسیدیم: \textbf{تکرار}. رایانه در
هیچ کاری به اندازهٔ تکرار خوب نیست؛ کاری را که برای ما کشنده و خسته‌کننده
است (جمع زدنِ هزار عدد، بررسیِ یک‌به‌یکِ میلیون‌ها حالت) بی‌حوصلگی و
بی‌خطا انجام می‌دهد. ابزارِ تکرار در برنامه‌نویسی \textbf{حلقه} نام دارد.

\section{تاهنگام: تکرار تا برقراریِ شرط}

بنیادی‌ترین حلقهٔ سلام \code{تاهنگام} است: بدنه‌اش را \emph{تا وقتی که} شرطش
درست است، بارها و بارها اجرا می‌کند. بیایید اعدادِ ۱ تا ۵ را چاپ کنیم:

\begin{salamfa}
کارکرد آغازین:
    گذرا i : صحیح = 1
    تاهنگام i <= 5:
        چاپ i
        i = i + 1
    پایان
    چاپ "تمام شد"
پایان
\end{salamfa}

\begin{progoutfa}
1
2
3
4
5
تمام شد
\end{progoutfa}

گردشِ کار چنین است: سلام شرط را می‌آزماید؛ اگر درست بود بدنه را اجرا
می‌کند و دوباره برمی‌گردد سراغِ شرط؛ و این چرخه ادامه دارد تا شرط نادرست
شود، که آن‌وقت حلقه تمام می‌شود و اجرا از بعدِ \code{پایان} پی گرفته می‌شود.
هر حلقهٔ خوش‌رفتار سه جزء دارد که در همین مثال پیداست:

\begin{enumerate}
  \item \textbf{آماده‌سازی:} \code{گذرا i = 1} پیش از حلقه؛
  \item \textbf{شرطِ ادامه:} \code{i <= 5}؛
  \item \textbf{گامِ پیشروی:} \code{i = i + 1} درونِ بدنه، که ما را
        قدم‌به‌قدم به پایان نزدیک می‌کند.
\end{enumerate}

\begin{warn}
گامِ پیشروی را جا بیندازید تا با ترسناک‌ترین اشتباهِ تازه‌کارها آشنا شوید:
\textbf{حلقهٔ بی‌پایان}. بدونِ \code{i = i + 1}، مقدارِ \code{i} همیشه ۱
می‌ماند، شرط همیشه درست است و برنامه تا ابد عدد ۱ چاپ می‌کند. اگر گرفتارش
شدید نترسید؛ در ترمینال کلیدهای \lr{Ctrl+C} را بزنید تا برنامه متوقف شود،
بعد گامِ پیشروی را پیدا و درست کنید.
\end{warn}

\section{الگوی انباشتن: جمع زدن در حلقه}

پرکاربردترین الگوی حلقه‌ها «انباشتن» است: متغیری بیرونِ حلقه بسازید و در هر
دور چیزی به آن بیفزایید. مجموعِ اعدادِ ۱ تا ۱۰۰ (مسئله‌ای که می‌گویند گاوس
در کودکی در یک چشم‌به‌هم‌زدن حل کرد) برای حلقه چیزی نیست:

\begin{salamfa}
کارکرد آغازین:
    گذرا مجموع : صحیح = 0
    گذرا i : صحیح = 1
    تاهنگام i <= 100:
        مجموع = مجموع + i
        i = i + 1
    پایان
    چاپ "مجموع ۱ تا ۱۰۰:", مجموع
پایان
\end{salamfa}

\begin{progoutfa}
مجموع ۱ تا ۱۰۰: 5050
\end{progoutfa}

همین الگو با کمی تغییر همه‌جا کار می‌کند: شمارنده بسازید و در هر دورِ
دلخواه یکی اضافه کنید تا «تعداد» به دست آید؛ در متغیرِ انباشته ضرب کنید تا
فاکتوریل به دست آید؛ و در فصل‌های بعد، بیشینه و کمینه را هم با همین الگو
پیدا می‌کنیم.

\section{چرخه: حلقهٔ شمارشی}

خیلی وقت‌ها از قبل می‌دانیم چند بار می‌خواهیم تکرار کنیم. برای این حالت
سلام حلقهٔ \code{چرخه} را دارد که کوتاه‌تر و خواناتر است:

\begin{salamfa}
کارکرد آغازین:
    چرخه 3:
        چاپ "هورا!"
    پایان
پایان
\end{salamfa}

\begin{progoutfa}
هورا!
هورا!
هورا!
\end{progoutfa}

اگر به شمارهٔ دور هم نیاز دارید، با واژهٔ \code{با} یک شمارنده بگیرید.
شمارنده از \emph{صفر} شروع می‌شود و تا یکی کمتر از عددِ چرخه پیش می‌رود:

\begin{salamfa}
کارکرد آغازین:
    چرخه 5 با i:
        چاپ "دور شماره", i
    پایان
پایان
\end{salamfa}

\begin{progoutfa}
دور شماره 0
دور شماره 1
دور شماره 2
دور شماره 3
دور شماره 4
\end{progoutfa}

\begin{note}
چرا از صفر؟ شمارش از صفر سنتِ ریشه‌دارِ دنیای برنامه‌نویسی است و در فصلِ
\ref{ch:arrays} که آرایه‌ها را می‌بینید، دلیلِ خوبش روشن می‌شود. فعلاً این
را به خاطر بسپارید: \code{چرخه 5 با i} پنج دور می‌زند و \code{i} از ۰ تا ۴
پیش می‌رود، نه ۱ تا ۵.
\end{note}

شکلِ سومِ \code{چرخه} بازه‌ای است: \code{چرخه 1 تا 10 با i:} شمارنده را
از ۱ تا ۱۰ (هر دو سر شامل) پیش می‌برد. جدول‌ضربِ ۷ در سه خط:

\begin{salamfa}
کارکرد آغازین:
    چرخه 1 تا 10 با i:
        چاپ 7, "*", i, "=", 7 * i
    پایان
پایان
\end{salamfa}

\section{بشکن و بگذر}

دو دستورِ کوچک، اختیارِ حلقه را کامل می‌کنند:

\begin{itemize}
  \item \code{بشکن} حلقه را همان لحظه به پایان می‌رساند و اجرا از بعدِ
        \code{پایان}ِ حلقه دنبال می‌شود.
  \item \code{بگذر} فقط \emph{همین دور} را رها می‌کند و یک‌راست سراغِ دورِ
        بعد می‌رود.
\end{itemize}

مثلاً در جست‌وجو، همین که گمشده پیدا شد دیگر ماندن در حلقه بی‌معنی است:

\begin{salamfa}
کارکرد آغازین:
    هدف : صحیح = 91
    گذرا n : صحیح = 2
    تاهنگام n <= 100:
        اگر n * n > هدف:
            بشکن
        پایان
        n = n + 1
    پایان
    چاپ "نخستین عددی که مربعش از ۹۱ بیشتر است:", n
پایان
\end{salamfa}

\begin{progoutfa}
نخستین عددی که مربعش از ۹۱ بیشتر است: 10
\end{progoutfa}

\section{حلقه درونِ حلقه}

بدنهٔ حلقه می‌تواند خودش حلقه داشته باشد. حلقهٔ بیرونی هر دورش، کلِ حلقهٔ
درونی را از اول تا آخر می‌چرخاند. برای چاپِ یک مثلثِ ستاره:

\begin{salamfa}
کارکرد آغازین:
    چرخه 1 تا 4 با سطر:
        چرخه سطر:
            بنویس "*"
        پایان
        چاپ ""
    پایان
پایان
\end{salamfa}

\begin{progout}
*
**
***
****
\end{progout}

سطرِ ۱ یک ستاره می‌گیرد، سطرِ ۲ دو ستاره و همین‌طور تا ۴. دستورِ
\code{چاپ ""} فقط به خطِ بعد می‌رود تا سطرِ بعدیِ مثلث از سرِ خط شروع
شود. حلقه‌های تودرتو را با اجرای دستی دنبال کنید تا کاملاً جا بیفتند؛
ارزشش را دارد، چون در جدول‌ها، شکل‌ها و مرتب‌سازیِ فصلِ
\ref{ch:classic-algorithms} دوباره سراغشان می‌رویم.

\section{کدام حلقه را انتخاب کنیم؟}

قاعدهٔ ساده: اگر \emph{تعدادِ} تکرار از پیش معلوم است، \code{چرخه}
خواناتر است؛ اگر پایانِ کار به \emph{شرطی} بسته است که وسطِ راه معلوم
می‌شود (تا وقتی خطا هست، تا وقتی عدد بزرگ‌تر نشده)، \code{تاهنگام} ابزارِ
درست است. هر دو هم‌ارزند؛ هر چه با یکی بشود نوشت با دیگری هم می‌شود، اما
خوانایی فرق می‌کند.

\section{تمرین‌ها}

\exercise با \code{تاهنگام} اعدادِ زوجِ ۲ تا ۲۰ را چاپ کنید. بعد همان را با
\code{چرخه} بازنویسی کنید.

\exercise شمارشِ معکوسِ موشک: از ۱۰ تا ۱ را چاپ کنید و در پایان
\code{"پرتاب!"}. (راهنمایی: شمارنده را بزرگ شروع کنید و در هر دور کم کنید.)

\exercise فاکتوریلِ ۱۰ را حساب کنید: حاصل‌ضربِ ۱ تا ۱۰. از الگوی انباشتن
استفاده کنید؛ فقط مقدارِ شروع را با دقت انتخاب کنید (چرا ۰ شروعِ خوبی
نیست؟).

\exercise برنامه‌ای بنویسید که رقم‌های عددِ \code{4725} را یکی‌یکی (از
یکان به بالا) چاپ کند. راهنمایی: ترکیبِ \code{\% 10} و \code{/ 10} در یک
حلقهٔ \code{تاهنگام عدد > 0}.

\exercise مثلثِ ستاره را وارونه چاپ کنید: سطرِ اول چهار ستاره، سطرِ آخر
یکی.

\exercise این حلقه چند بار «سلام» چاپ می‌کند؟ اول روی کاغذ، بعد با اجرا:
\begin{salamfa}
کارکرد آغازین:
    گذرا i : صحیح = 10
    تاهنگام i > 0:
        چاپ "سلام"
        i = i - 3
    پایان
پایان
\end{salamfa}
